\documentclass[pdf, 10pt]{beamer}

% =========================================================
% Theme and Colors
% =========================================================
\usetheme{Madrid}
\usecolortheme{whale}

% =========================================================
% Packages
% =========================================================
\usepackage[utf8]{inputenc}
\usepackage[T1]{fontenc}
\usepackage{graphicx}
\usepackage{booktabs}
\usepackage{bookmark}

% =========================================================
% Presentation Info
% =========================================================
\title[AI Mobile Mast Handover]{AI-Based Mobile Perception System}
\subtitle{Project Handover and Technical Roadmap}
\author[Batch 2025-26]{Asfak, Sahil, Shwetha, Sanjay, Fatih}
\institute[THI]{Technische Hochschule Ingolstadt}
\date[Feb 2026]{February 6, 2026}

\begin{document}

% ---------------------------------------------------------
% Title Slide
% ---------------------------------------------------------
\begin{frame}
    \titlepage
\end{frame}

% ---------------------------------------------------------
% Slide 1: Introduction
% ---------------------------------------------------------
\begin{frame}{Project Objective}
    \begin{itemize}
        \item \textbf{Goal:} Develop a robust, mobile surveillance system using a mast platform.
        \item \textbf{Platform:} Distributed processing across 3x NVIDIA Jetson Orin devices.
        \item \textbf{Sensing:} Multi-modal fusion of 128-beam LiDAR and 360-degree Camera array.
        \item \textbf{Output:} Real-time 3D object detection, tracking, and global visualization.
    \end{itemize}
\end{frame}

% ---------------------------------------------------------
% Slide 2: Distributed Architecture
% ---------------------------------------------------------
\begin{frame}{System Architecture (Network & Hardware)}
    \begin{columns}
        \column{0.5\textwidth}
        \textbf{Orin 1: Sensor Gateway}
        \begin{itemize}
            \item Raw LiDAR (RS16) acquisition.
            \item Camera RTSP stream handling.
            \item Preprocessing & Clustering.
        \end{itemize}
        
        \textbf{Orin 2: Perception Engine}
        \begin{itemize}
            \item VPI GPU-accelerated undistortion.
            \item YOLOv11s object detection.
            \item 2D-3D projection (Cam-LiDAR).
        \end{itemize}

        \column{0.5\textwidth}
        \textbf{Orin 3: Global Tracker}
        \begin{itemize}
            \item Kalman Filter-based tracking.
            \item Appearance Re-ID (DeepSORT).
            \item Centralized Mission Control GUI.
        \end{itemize}
    \end{columns}
\end{frame}

% ---------------------------------------------------------
% Slide 3: Key Technology: Time Synchronization
% ---------------------------------------------------------
\begin{frame}{Ensuring Temporal Alignment (PTP & Chrony)}
    \begin{itemize}
        \item \textbf{The Problem:} Ghosting artifacts in fusion due to sensor clock drift.
        \item \textbf{The Solution:} Two-layer network synchronization.
        \item \textbf{Layer 1 (Sensor):} Precision Time Protocol (IEEE 1588) for sub-microsecond LiDAR-Camera sync.
        \item \textbf{Layer 2 (System):} Chrony NTP for consistent ROS 2 message timestamps across Orins.
        \item \textbf{Result:} Fusion spatial error reduced to $<15$cm at urban traffic speeds.
    \end{itemize}
\end{frame}

% ---------------------------------------------------------
% Slide 4: Results and Performance
% ---------------------------------------------------------
\begin{frame}{Key Performance Metrics}
    \begin{table}
        \centering
        \begin{tabular}{l c c}
            \toprule
            \textbf{Component} & \textbf{Latency} & \textbf{Update Rate} \\
            \midrule
            LiDAR Clustering & 12ms & 10 Hz \\
            YOLO Detection (FP16) & 22ms & 30 Hz \\
            Tracker & 8ms & 30 Hz \\
            \midrule
            \textbf{End-to-End System} & \textbf{$\approx$49ms} & \textbf{20 Hz} \\
            \bottomrule
        \end{tabular}
    \end{table}
    \begin{itemize}
        \item mAP@0.5 Detection Accuracy: \textbf{0.87} (Vehicles, Pedestrians, Cyclists).
        \item Stable tracking IDs for $>95\%$ of valid targets.
    \end{itemize}
\end{frame}

% ---------------------------------------------------------
% Slide 5: Handover: How to Start
% ---------------------------------------------------------
\begin{frame}{Handover: How to Start}
    \begin{block}{Static IP Scheme (Subnets 192.168.6.x & 192.168.51.x)}
        \begin{itemize}
            \item Central PC: \texttt{.105} | Orin 1: \texttt{.106} | Orin 2: \texttt{.107} | Orin 3: \texttt{.108}
            \item LiDAR: \texttt{51.132} | Cameras: \texttt{51.133 - .139}
        \end{itemize}
    \end{block}
    \begin{block}{Quick Launch Steps}
        \begin{enumerate}
            \item Start PTP Master (Orin 1/Switch).
            \item Establish SSH connection via Mission Control GUI.
            \item Trigger "Connect All" $\rightarrow$ "Sequential Launch".
            \item Verify markers in RViz.
        \end{enumerate}
    \end{block}
\end{frame}

% ---------------------------------------------------------
% Slide 6: Future Roadmap (What to do next)
% ---------------------------------------------------------
\begin{frame}{Technical Roadmap & Future Tasks}
    \begin{enumerate}
        \item \textbf{Appearance Modeling:} Enhance Re-ID with Attention-based CNNs for crowded scenarios.
        \item \textbf{Adverse Weather:} Integrate Thermal (FLIR) cameras to handle night/fog detection.
        \item \textbf{Compute Optimization:} Utilize NVIDIA DeepStream for zero-copy GPU memory handling between VPI and YOLO.
        \item \textbf{Mobile Communication:} Test system stability over high-latency 5G/Rajant wireless links.
    \end{enumerate}
\end{frame}

% ---------------------------------------------------------
% Final Slide
% ---------------------------------------------------------
\begin{frame}
    \centering
    \Huge Thank You! \\
    \vspace{1cm}
    \Large Questions? \\
    \vspace{1cm}
    \small Source Core: \url{https://github.com/Asfak3566/AI-Mobile-Perception-System}
\end{frame}

\end{document}
